\documentclass[11pt,a4paper]{article}
\usepackage[T1]{fontenc}
\usepackage{lmodern,amsmath,amssymb,mathtools,booktabs,tabularx,microtype,xcolor,hyperref}
\usepackage[margin=26mm]{geometry}
\usepackage{enumitem,fancyhdr}
\definecolor{ink}{HTML}{244B68}
\hypersetup{colorlinks=true,linkcolor=ink,urlcolor=ink,pdftitle={BundleFuse: A Theory of Completing Bundles}}
\pagestyle{fancy}\fancyhf{}\fancyhead[L]{\small BundleFuse: completing bundles}\fancyhead[R]{\small Theory}\fancyfoot[C]{\thepage}
\setlength{\headheight}{14pt}\setlength{\parindent}{0pt}\setlength{\parskip}{6pt}\setlength{\emergencystretch}{2em}
\newcommand{\R}{\mathbb R}\newcommand{\I}{\mathcal I}\newcommand{\diag}{\operatorname{diag}}
\newcommand{\off}{\operatorname{off}}\newcommand{\pos}[1]{\left[#1\right]_+}
\newcommand{\normrow}{\mathcal N}\newcommand{\sd}{\operatorname{sd}}
\numberwithin{equation}{section}
\title{\textcolor{ink}{BundleFuse}\\[5pt]\large A Theory of Completing Bundles}
\author{}\date{}
\begin{document}\maketitle
\begin{abstract}
How can we suggest an item that fits a partly completed collection? This chapter develops
an answer from first principles. It begins with past collections and the items they contain,
then builds several kinds of evidence: shared history, support from groups of known items,
item descriptions, and categories. It explains how those pieces are put on a common scale,
combined, and adjusted when an item has little history. A small numerical example follows
the calculation through to a ranked recommendation.
\end{abstract}
\tableofcontents\clearpage
\section{The big picture: from a partial bundle to suggestions}
A \emph{bundle} is a collection of items that belong together. A bundle could be an outfit,
a shopping basket, or a music playlist. Suppose we are shown a few items from a bundle and
asked to suggest the missing items. The items already shown are called the \emph{known items}.
A possible missing item is called a \emph{candidate}.

A good candidate should fit the known items as a group. History helps: items previously
found in the same bundles often fit together. Descriptions and categories help too: a new
item may have no purchase or playlist history but may still resemble useful items. These
clues need not agree. BundleFuse assigns a score to every candidate, combines the clues,
and returns the highest-scoring items that are not already known.

\subsection{One question, five steps}
Imagine a store with five items, numbered 0 to 4. Someone has chosen items 0 and 1. We
want to decide whether to suggest item 2, 3, or 4. Suppose two earlier baskets were
$\{0,1,2\}$ and $\{0,3\}$. Item 2 appeared with both chosen items, item 3 appeared
with one of them, and item 4 has no basket history. The answer is not automatic: item 4
may still have a useful description, while an item seen often may simply be popular.

BundleFuse addresses this decision in five connected steps:
\begin{enumerate}[leftmargin=*]
\item \textbf{Gather evidence.} Look at past bundles and, when available, item descriptions,
      categories, and other relationships. Each kind of evidence gives each candidate a score.
\item \textbf{Make the scores comparable.} A count and a description similarity have
      different units. Rescale each kind of evidence across the available items.
\item \textbf{Combine the scores.} Give each kind of evidence an importance, and allow
      extra credit when several known items support a candidate together.
\item \textbf{Account for limited history.} Blend the combined score with an estimate
      based on descriptions and categories, giving that estimate more influence for an item
      with little history.
\item \textbf{Rank and check.} Remove the items already chosen, rank the remaining items,
      and assess the suggestions on examples whose hidden items are known.
\end{enumerate}
The following sections explain each step in this order. We start by putting past bundles
into a table, because every score based on shared history needs that table. We return to
the five-item example throughout, and finish its numerical calculation in
Section~\ref{example}. The equations make the rules precise; the paragraphs around them
explain what each rule contributes to the decision.

\subsection{The symbols we will use}
\begin{center}\small
\begin{tabularx}{\linewidth}{@{}p{.20\linewidth}X@{}}\toprule
Symbol & Meaning \\\midrule
$\I$ & Set of all available items; $n=|\I|$ is the number of items. \\
$b_e$ & Past bundle number $e$; there are $m$ past bundles. \\
$S$ & Set of known items for the bundle we want to complete. \\
$i,s,t$ & Item labels: $i$ is usually a candidate; $s,t$ are known items. \\
$B$ & A table with one row per past bundle and one column per item. \\
$f_i$ & Number of past bundles containing item $i$. \\
$d_e$ & Number of items in past bundle $e$. \\
$r_i^h(S)$ & Raw evidence for candidate $i$ from source $h$, given $S$. \\
$z_i^h(S)$ & The same evidence after it has been put on a common scale. \\
$w_h$ & Nonnegative importance assigned to source $h$. \\
$Y_i(S)$ & Combined evidence from the available sources. \\
$P_i(S)$ & A second, content-oriented estimate used when history is limited. \\
$\widehat\theta_i(S)$ & Final score for candidate $i$. \\
$K$ & Number of suggestions requested. \\
\bottomrule\end{tabularx}\end{center}
We use $[x]_+=\max(x,0)$, so $[-2]_+=0$ and $[3]_+=3$. A small positive number $\varepsilon$
in a denominator prevents division by zero. When a number is stated as a usual setting, it
is a choice of the method rather than a mathematical necessity.

\section{Representing past bundles}
Define the entry in row $e$ and column $i$ of $B$ as
\begin{equation}
 B_{ei}=\begin{cases}1,&i\in b_e,\\0,&i\notin b_e.\end{cases}
\end{equation}
The number of past bundles containing item $i$, and the size of bundle $e$, are
\begin{equation}
 f_i=\sum_{e=1}^{m}B_{ei},\qquad d_e=\sum_{i\in\I}B_{ei}.
\end{equation}
For example, if the past bundles are \{tea, milk, sugar\} and \{tea, biscuits\}, then
tea occurs twice, and the first bundle contains three items. If two different past bundles
contain exactly the same items, they count as two observations.

For the current request, let $x_i=1$ when $i\in S$ and $x_i=0$ otherwise. The vector $x$
therefore marks the known items. For our example, the table has two rows: the first marks items 0, 1, and 2, and the second
marks 0 and 3. The indicator marks items 0 and 1. This gives us the basic information
needed for the first question: how much does each past bundle support each candidate?

\section{Evidence from shared history}
\subsection{Simple counts}
For distinct items $i$ and $j$, let $C_{ij}$ be the number of past bundles containing both.
Set $C_{ii}=0$. In matrix form,
\begin{equation}
 C=\operatorname{off}(B^\top B),
\end{equation}
where $B^\top$ turns rows into columns and $\operatorname{off}$ replaces diagonal entries
by zero. A simple score for $i$ is $\sum_{s\in S} C_{is}$: add the number of times $i$
appeared with each known item. In the example, item 2 appeared with both known items,
while item 3 appeared with item 0 only. This is easy to understand, but very common items can
receive large scores merely because they appear in many bundles.

\subsection{Giving each past bundle a suitable weight}
Let $\alpha_d\ge0$ and $\gamma_v\ge0$ control two adjustments. A usual choice is
$\alpha_d=\gamma_v=\tfrac12$. Define the weight of a known item $s$ and the weight of a
past bundle $e$ by
\begin{equation}
 a_s=(f_s+5)^{-\alpha_d},\qquad v_e=(d_e+1)^{-\gamma_v}.
\end{equation}
The first reduces the influence of a very common known item. The second reduces the
influence of a very large past bundle. The constants 5 and 1 keep the weights finite and
soften the adjustment for small counts.

The weighted overlap between past bundle $e$ and the known items is
\begin{equation}\label{overlap}
 o_e(S)=v_e\sum_{s\in S}B_{es}a_s.
\end{equation}
The first history-based score of candidate $i$ is the sum of these overlaps over past
bundles containing it:
\begin{equation}\label{pair}
 p_i(S)=\sum_{e=1}^{m}B_{ei}\,o_e(S).
\end{equation}
A candidate is supported when it occurs in past bundles that also contain known items.
Equation~\eqref{pair} adds the contributions of known items; it does not demand that every
known item be present in the same past bundle.

\subsection{Extra support from several known items}
Let $k_e(S)=\sum_{s\in S}B_{es}$ be the number of known items in past bundle $e$.
Let $\beta\ge0$ control the extra credit for group overlap; a usual value is $1.5$.
Define
\begin{equation}\label{surplus}
 \delta_i(S)=\sum_{e=1}^{m}B_{ei}\,o_e(S)\,\beta[k_e(S)-1]_+.
\end{equation}
If a past bundle contains one known item, its extra contribution is zero. If it contains
two, it receives extra credit; three receive more. The strengthened history score is
$p_i(S)+\delta_i(S)$. With only one known item, $\delta_i(S)=0$ for every candidate.
In our example, this extra support goes to item 2 through the first basket. The second
basket cannot contribute extra group support to item 3, because it contains only one known
item. We now have a direct score and a group-strengthened score. The next step asks whether
connections through other items can provide additional evidence.

\subsection{Following connections one step farther}
A candidate can be related to the known items through other items even if the direct
connection is weak. To describe this, let $D$ be the diagonal table with entries
$D_{ii}=(f_i+5)^{-\alpha_d}$, and let $E$ be the diagonal table with entries
$E_{ee}=(d_e+1)^{-1}$. Define the item-to-item table
\begin{equation}\label{L}
 L=\operatorname{off}(DB^\top EBD).
\end{equation}
For any vector $a$ of item scores, $La$ spreads those scores across items that have appeared
in the same past bundles, while the zero diagonal prevents an item from sending its score
straight back to itself. The two longer-range scores are $Lp$ and $L(p+\delta)$.
For shopping data with additional user and session histories, a different table can take
the place of $L$; Section~\ref{context} explains that variation.

\section{When pairs occur together more often than expected}
The preceding scores use observed overlaps, adjusted for the frequency of known items and
the size of past bundles. There is another question: did two items occur together more often
than we would expect from how common they are individually? Answering it gives us a
second view of the same history.

A raw shared-history count can be large simply because both items are common. Let
$g_i=\sum_j C_{ij}+\varepsilon$ be the total shared-history count of item $i$, and let
$Z_C=\sum_{i,j}C_{ij}$ be the total of the table. For a pair observed at least once, define
\begin{equation}\label{pmi}
 M_{ij}=\left[\log\left(\frac{C_{ij}Z_C}{g_i g_j}+10^{-12}\right)\right]_+.
\end{equation}
For a pair never observed together, set $M_{ij}=0$. This number is large when the pair's
count is strong relative to its overall counts, and zero when the expression is not positive.
A candidate's score from this source is $\sum_{s\in S}M_{is}$.

A second source makes a shorter description of the rows of $M$. Choose a dimension $r$;
find a rank-$r$ approximation to $M$ by a method called \emph{singular value decomposition},
which finds directions that summarize much of a table's variation. Each item then has a
short vector $u_i\in\R^r$; normalize each nonzero vector to length approximately one.
The score is the dot product of the candidate vector with the mean known-item vector:
\begin{equation}\label{meanfeature}
 r_i^{\mathrm{history}}(S)=u_i\cdot\left(\frac1{|S|}\sum_{s\in S}u_s\right).
\end{equation}
A dot product adds the products of corresponding vector entries. Here it measures how
closely the candidate's summary of history agrees with the average known item.

\section{Evidence from descriptions and categories}\label{content}
History can tell us little about an item that rarely appears, and nothing about an item
with no appearances. In our example, item 4 would receive no direct shared-history support.
Its description or category may still tell us something useful. We now develop those sources,
which will later give such items a way to compete with well-observed candidates.

\subsection{Comparing item descriptions}
Suppose item $i$ has a vector $c_i$ describing its text, image, or audio. These vectors
can come from measurements or from an existing representation of the item. Normalize the
vectors and compare a candidate with the mean description of the known items:
\begin{equation}
 r_i^{\mathrm{description}}(S)=c_i\cdot\left(\frac1{|S|}\sum_{s\in S}c_s\right).
\end{equation}
A dot product near one means that normalized vectors point in similar directions; near
zero indicates little agreement in that representation. Different types of description
can produce separate scores and receive different weights.

For text, one possible description counts words and two-word phrases, gives less weight
to terms appearing in many item descriptions, and then shortens the resulting table by
singular value decomposition. The resulting coordinates summarize descriptions; they do
not require that a candidate has appeared in a past bundle.

\subsection{Separating category from the rest of a description}
Items in the same category may share broad features that say little about which particular
items fit together. Let $c_i$ now denote a vector from one description type, and let
$g(i)$ be the group of item $i$. Write $\bar c_g$ for the average vector of all items in
group $g$. The remaining description after removing the group average is
\begin{equation}
 d_i=\frac{c_i-\bar c_{g(i)}}{\lVert c_i-\bar c_{g(i)}\rVert_2+\varepsilon}.
\end{equation}
Here $\lVert\cdot\rVert_2$ is ordinary vector length. For an outfit, the group might be a
clothing category; for audio, it may be a group of items with similar sound. The vectors
$d_i$ describe how items differ within their broad groups. They can be scored by the same
mean-and-dot-product rule as equation~\eqref{meanfeature}.

\subsection{Joining descriptions with history}
Let $D_c$ be the table whose row $i$ is $d_i$. Apply $L$ from equation~\eqref{L} to its
columns, and normalize each resulting row; call that normalized table $H$. For a chosen
weight $\lambda\ge0$, form and row-normalize
\begin{equation}
 F=\operatorname{normalize\ rows}(D_c+\lambda H).
\end{equation}
A row of $F$ blends an item's own description with descriptions of historically connected
items. Comparing its rows using equation~\eqref{meanfeature} supplies another score.
A usual value is $\lambda=0.5$.

\subsection{Do categories complement each other?}
Let $g(i)$ be the category of item $i$. Define $A_{ec}$ as the number of items of category
$c$ in past bundle $e$, and let $a_c=\sum_e A_{ec}+\varepsilon$. Define
$H=\operatorname{off}(A^\top A)$ and $A_+=\sum_c a_c$. For observed category pairs,
\begin{equation}
 Q_{cd}=\left[\log\left(\frac{H_{cd}A_+}{a_c a_d}+10^{-12}\right)\right]_+,
\end{equation}
and set unobserved entries to zero. The category score for candidate $i$ is
\begin{equation}
 r_i^{\mathrm{category}}(S)=\sum_{s\in S}Q_{g(s),g(i)}.
\end{equation}
For example, if shirts and trousers frequently occur together, that relationship can help
score trousers when a shirt is known. In some applications, an additional penalty discourages
repeating a category that seldom appears twice in the same past bundle. In others, shared
artist or album membership adds a separate score. The appropriate category evidence
depends on what a bundle represents.

\section{Additional relationships in shopping data}\label{context}
The past-bundle table and item descriptions are not always the only available evidence.
In shopping data, for example, several baskets may belong to the same household. These
relationships supply another view of how items are connected. They are optional: the
reasoning up to this point works without them.

Some data contain households, users, sessions, or an additional record of user--item
interactions. These can provide extra evidence without changing the purpose of the method.
For example, create a table $Z$ whose rows are households and whose columns are items,
with each entry counting appearances in a household's past training baskets. Divide each
household row by the square root of one plus its total count. The product of the resulting
table's transpose and itself measures how often items are associated through households;
its diagonal is removed. A candidate can then receive evidence when it shares household
history with the known items.

For some commerce data, separate user and session tables give item-pair tables $C_u$ and
$C_s$, and an additional interaction record gives a weighted pair table $C_p$. Define
$D_c=\diag((f_i+1)^{-1/2})$. One longer-range table is
\begin{equation}
 L_c=D_c(C+0.5C_s+0.5C_u+0.2C_p)D_c.
\end{equation}
It spreads evidence through several kinds of shared history. The particular coefficients
are choices that define this version of the method. Their presence means that a result for
this setting should be understood in the light of the extra information it uses.

\section{Agreement among known items}\label{agreement}
So far, we have considered whether each candidate has historical, descriptive, or category
support. We can also ask how that support is distributed across the known items. Is it
coming from just one known item, or do several known items agree? This question matters
in our example because item 2 is connected to both 0 and 1.

A candidate is especially interesting when several known items support it together. To
measure this, let $q_{is}\ge0$ be the support that known item $s$ gives candidate $i$.
For a concrete definition, let $W_{is}=0$ when $i=s$; otherwise let
$W_{is}=\sum_e B_{ei}v_e B_{es}a_s$, using the weights defined above. Let $d_i$ be a
normalized description vector and $\lambda_d\ge0$ a chosen description weight. Then
\begin{equation}
 q_{is}=W_{is}+\lambda_d[d_i\cdot d_s]_+,\qquad
 A_i=\sum_{s\in S}q_{is}.
\end{equation}
If description vectors are unavailable, the second term is omitted.

Known items may be nearly duplicates of one another. Let $G_{st}$ measure similarity
between known items $s$ and $t$. When both category labels $g(\cdot)$ and description
vectors are available, one choice is
\begin{equation}
 G_{st}=\tfrac12\mathbf1[g(s)=g(t)]+\tfrac12[d_s\cdot d_t]_+.
\end{equation}
With only one kind of information, use that part alone. For distinct known items, choose a coefficient $\eta$ and define
\begin{equation}
 K_{st}=\begin{cases}[1+\eta G_{st}]_+,&\eta<0,\\1+\eta G_{st},&\eta\ge0,\end{cases}
 \qquad K_{ss}=0.
\end{equation}
A negative $\eta$ reduces the extra credit for similar known items; a positive one can
increase it. Let $\tau>0$ be a small stabilizing number, commonly $0.05$. Define the
agreement score
\begin{equation}\label{agreementeq}
 r_i^{\mathrm{agreement}}(S)=
 \frac{\displaystyle\sum_{s\in S}\sum_{t\in S}q_{is}K_{st}q_{it}}{A_i+\tau}.
\end{equation}
Both orders of a distinct pair are included in the double sum. If fewer than two known
items are given, the score is zero. This score rewards support from pairs of known items;
it does not require the candidate to match every item simultaneously.

\section{Putting different evidence on the same scale}\label{scale}
We now have several possible scores for the same candidate. To make one decision, we need
to combine them. Before doing so, we must address a simple problem: counts, dot products,
and category scores have different units. A count of 10 is not automatically stronger than
a similarity of 0.8.

Counts, dot products, and category scores have different units. For each evidence source
$h$ and a fixed set of known items $S$, let $r_i^h(S)$ be its raw score for item $i$.
Its mean across all $n$ items is
\begin{equation}
 \mu_h(S)=\frac1n\sum_{i\in\I}r_i^h(S).
\end{equation}
Let $\sigma_h(S)$ be the corresponding standard deviation: a measure of how much the
scores vary around their mean. The common-scale score is
\begin{equation}\label{standard}
 z_i^h(S)=\frac{r_i^h(S)-\mu_h(S)}{\sigma_h(S)+\varepsilon}.
\end{equation}
A positive value means above-average evidence for that source; a negative value means
below average. If all items receive the same raw score, the common-scale scores are zero.
The mean and spread are calculated for each request across the whole item set, before the
known items are removed from the final list. Some settings use a sample standard deviation
(dividing the sum of squared deviations by $n-1$); others use a population standard deviation
(dividing by $n$). Keeping this choice fixed matters for exact comparisons.

The strengthened history score and the simple history score are each standardized
separately. Therefore, the standardized value of $p+\delta$ is generally not the sum of
the separately standardized values of $p$ and $\delta$.

\section{Combining the evidence}\label{combine}
The preceding step expresses each kind of evidence relative to its own average and spread.
We can now decide how much importance to give each kind and combine them into one score.
This yields an answer based on the available observations; the next section will adjust
that answer for items whose history is scarce.

Let $\mathcal H$ be the set of evidence sources under consideration, and let $w_h\ge0$
be the importance of source $h$. For each source, let $z_i^{h,0}$ be its score with the
extra group contribution removed; for sources without such a contribution, let
$z_i^{h,0}=z_i^h$. Define
\begin{equation}
 Y_i^0=\sum_{h\in\mathcal H}w_h z_i^{h,0},\qquad
 Y_i^1=\sum_{h\in\mathcal H}w_h z_i^h.
\end{equation}
Choose $a\in[0,1]$ for the amount of group evidence, and $b\ge0$ for the amount of
agreement evidence from equation~\eqref{agreementeq}. The combined observed score is
\begin{equation}\label{observed}
 Y_i=(1-a)Y_i^0+aY_i^1+bz_i^{\mathrm{agreement}}.
\end{equation}
When $a=0$, the calculation uses the simpler history scores. When $a=1$, it uses their
strengthened counterparts. The agreement term is an additional piece of evidence.

The weights are selected using held-out examples for which the hidden items are known.
One can try a short list of nonnegative values for one weight at a time, retain a change
when it improves a chosen ranking measure, and repeat. A small tolerance can remove a
source that adds little. The amounts $a$ and $b$ are then chosen from small lists. This is
a search over settings informed by examples; it does not require a neural network.
The held-out examples used for selecting settings are kept separate from the final test
examples used to report performance.

\section{Making a sensible choice when history is limited}\label{shrink}
The combined score answers what the available evidence suggests. But an item seen once and
an item seen hundreds of times do not have equally rich histories. In our example, item 4
has no past-bundle history at all. We therefore bring in a second estimate that can rely
on information available even for that item.

A rare item has less historical evidence than a frequent one. BundleFuse therefore also
forms a second estimate from selected sources such as descriptions and categories. Let
$P_i$ be their weighted, common-scale sum. Depending on the setting, it can also include
the same agreement or popularity contributions used in the combined score. Let $s_0\ge0$
be a scale chosen using the held-out selection examples, and define the second estimate
$\mu_i=s_0P_i$.

Let $\kappa\ge0$ control how strongly the method relies on this second estimate. If
$\kappa=0$, the final score is simply $Y_i$. For $\kappa>0$, define
\begin{equation}\label{final}
 u_i=\frac{\kappa}{f_i+\kappa},\qquad
 \widehat\theta_i=(1-u_i)Y_i+u_i\mu_i.
\end{equation}
The weights in this average add to one. An item with $f_i=0$ receives the second estimate
entirely. As $f_i$ increases, $u_i$ decreases, and the observed score $Y_i$ has more
influence. For example, with $\kappa=2$, an item seen in two past bundles has $u_i=1/2$;
an item seen in 18 has $u_i=1/10$. If the method selects $\kappa=0$, this adjustment is
disabled for every item.

\subsection{A probability interpretation}
The weighted average above has a useful interpretation, provided we state its assumptions.
Suppose an unobserved true suitability $\theta_i$ follows a normal distribution with mean
$\mu_i$ and variance $\sigma_0^2$. Suppose also that the observed score $Y_i$, conditional
on $\theta_i$, has mean $\theta_i$ and variance $\sigma_e^2/f_i$ for $f_i>0$. Assume
$\sigma_0^2$ and $\sigma_e^2$ are positive. Then the conditional mean of $\theta_i$ given
$Y_i$ is equation~\eqref{final} with
\begin{equation}
 \kappa=\frac{\sigma_e^2}{\sigma_0^2}.
\end{equation}
The assumption says that more appearances reduce uncertainty in the observed evidence.
At $f_i=0$, the formula uses the prior-only limit. In practice, the final $\kappa$ can be
chosen from several values using held-out examples, so it need not equal a single fitted
variance ratio. The final scores are ranking scores, not probabilities that items belong
in the bundle; the assumptions do not by themselves establish that they are calibrated.

\subsection{Estimating a useful range for the adjustment}
Let there be $v$ held-out selection requests, and write $Y_{ij}$ and $\mu_{ij}$ for item
$i$ on request $j$. For each item with $f_i>0$, calculate its average squared difference
between the two estimates and its inverse history count:
\begin{equation}
 e_i=\frac1v\sum_{j=1}^{v}(Y_{ij}-\mu_{ij})^2,\qquad x_i=\frac1{f_i}.
\end{equation}
Fit a straight line to points $(x_i,e_i)$ across items. Its intercept suggests the prior
variance and its slope suggests the extra observation variance. If there are fewer than two
items with history, or all positive history counts are effectively equal, the two parts
cannot be distinguished from these points and the adjustment is disabled. Otherwise,
small positive lower bounds keep the estimated variances usable, and their ratio supplies
one candidate value of $\kappa$. Several fixed values and this fitted value can be compared
on the held-out selection examples. These squared differences compare two score estimates;
they are not squared prediction errors against the hidden-item labels.

\section{Putting the pieces together: the five-item example}\label{example}
We have now defined all five steps in the map at the beginning. Let us return to the same
five items and follow a candidate from past baskets to final score. This makes clear where
the group evidence and the adjustment for limited history enter.

Consider five available items numbered 0 through 4. The two past bundles are
$b_1=\{0,1,2\}$ and $b_2=\{0,3\}$; the known items are $S=\{0,1\}$. Item 4 has no past
bundle history. The history counts are $(f_0,f_1,f_2,f_3,f_4)=(2,1,1,1,0)$.

Use $\alpha_d=\gamma_v=\tfrac12$ and $\beta=1.5$. Define $a=1/\sqrt7$ and
$b=1/\sqrt6$. The two past-bundle overlaps from equation~\eqref{overlap} are
\begin{equation}
 o_1=\frac{a+b}{2},\qquad o_2=\frac{a}{\sqrt3}.
\end{equation}
The first bundle contains both known items; the second contains just item 0. Equations
\eqref{pair} and \eqref{surplus} give, in item order,
\begin{equation}
 p=(o_1+o_2,o_1,o_1,o_2,0),\quad
 \delta=(1.5o_1,1.5o_1,1.5o_1,0,0).
\end{equation}
Item 2 receives extra group support because it appeared with both 0 and 1. Item 3 receives
history support from item 0 but no group extra. Item 4 has no history support.

For a concrete illustration, suppose the selected settings use only the standardized
$p+\delta$ as $Y$, and the second estimate is half of a standardized description score.
Suppose those description scores after standardization are approximately
$(-1.265,-0.632,0,0.632,1.265)$, and use $\kappa=1$. The final scores for the three
candidates are approximately
\begin{center}\begin{tabular}{@{}lrrr@{}}\toprule
Candidate & History count $f_i$ & Weight $u_i$ & Final score \\\midrule
2 & 1 & $1/2$ & 0.288 \\
3 & 1 & $1/2$ & $-0.273$ \\
4 & 0 & $1$ & 0.632 \\\bottomrule\end{tabular}\end{center}
The top two suggestions are therefore items 4 and 2. Items 0 and 1 were included when
standardizing the evidence but are removed from the suggestion list because they were
already known. The values chosen here illustrate the calculation; different data and
held-out selection examples can lead to different weights and rankings.

\section{Choosing suggestions and measuring their quality}
The example produced a score for each candidate, but a practical answer needs a short list.
After forming that list, we need a separate way to check whether it recovers items that
were actually hidden. This final step connects the score to a measurable completion task.

Once every item has a final score $\widehat\theta_i$, remove all items already in $S$.
Sort the remaining items by decreasing score and return the first $K$, or all remaining items if fewer than $K$ are available. When two scores are
exactly equal, order the tied items by their item labels, from smaller to larger. This makes
the answer repeatable.

For evaluation, assume $K$ does not exceed the number of available candidates. Take a complete bundle and hide some items. Let $G$ be its nonempty set of
hidden items, and let $R_K$ be the top-$K$ suggestions from the revealed items. Recall at
$K$ is
\begin{equation}
 \operatorname{Recall}@K=\frac{|R_K\cap G|}{|G|}.
\end{equation}
It measures what fraction of the hidden items were recovered. It does not distinguish a
hit at rank 1 from a hit at rank $K$.

To value earlier hits more, let $h_r=1$ when the item at rank $r$ is hidden ground truth,
and $h_r=0$ otherwise. The discounted gain and its best possible value are
\begin{align}
 \operatorname{DCG}@K&=\sum_{r=1}^{K}\frac{h_r}{\log_2(r+1)},\\
 \operatorname{IDCG}@K&=\sum_{r=1}^{\min(K,|G|)}\frac1{\log_2(r+1)}.
\end{align}
The normalized discounted gain is $\operatorname{NDCG}@K=\operatorname{DCG}@K/
\operatorname{IDCG}@K$. A hit near the top contributes more. Scores for multiple bundles
are averaged, giving every evaluated bundle the same weight. Selection examples choose
settings; a separate test set measures their performance after those settings are fixed.

\section{What the theory does and does not establish}
The method combines several interpretable pieces of evidence, and the exposure-based
average offers a clear way to rely more on descriptions when history is scarce. The
construction alone does not prove that the pieces improve a particular dataset or that the
normal-distribution interpretation is accurate. Those questions require comparisons on
matched data, splits, inputs, and evaluation rules. Scores and rankings describe the method's
preferences under the available evidence; they do not establish causation or guarantee that
all recommended items are appropriate for every user.
\end{document}
